\documentclass[11pt]{article}
\usepackage[margin=1in]{geometry}
\usepackage{amsmath,amssymb,amsthm,mathtools}
\usepackage{bm}
\usepackage{booktabs}
\usepackage{array}
\usepackage{hyperref}
\usepackage{enumitem}

\hypersetup{
    colorlinks=true,
    linkcolor=blue,
    citecolor=blue,
    urlcolor=blue
}

\title{Report on the Two-Stage Offset-Free Target Selector}
\author{}
\date{\today}

\newcommand{\R}{\mathbb{R}}
\newcommand{\xaug}{x^{\mathrm{aug}}}
\newcommand{\xs}{x_s}
\newcommand{\us}{u_s}
\newcommand{\ds}{d_s}
\newcommand{\ysp}{y_{\mathrm{sp}}}
\newcommand{\yexpr}{y_s}
\newcommand{\yctrl}{y_c}
\newcommand{\dhat}{\hat d}
\newcommand{\xhataug}{\hat x^{\mathrm{aug}}}
\newcommand{\norminf}[1]{\left\lVert #1 \right\rVert_{\infty}}

\begin{document}
\maketitle

\section{Purpose}
This report documents the target selector currently implemented in
\texttt{Lyapunov/target\_selector.py} through the function
\texttt{compute\_ss\_target\_refined\_rawlings(...)}. The selector is designed
for the standard Lyapunov tracking MPC path in this repository and computes a
steady-state equilibrium target consistent with the current disturbance estimate.

The implementation follows a two-stage logic:
\begin{enumerate}[label=\arabic*.]
    \item solve an exact offset-free steady-state target problem, and
    \item if exact tracking is not admissible, solve a closest-feasible
    steady-state target problem.
\end{enumerate}

This structure is motivated by the standard disturbance-augmented offset-free MPC
formulation of Pannocchia and Rawlings \cite{PannocchiaRawlings2003},
Pannocchia and Bemporad \cite{PannocchiaBemporad2007}, and related tutorial and
review material \cite{MaederMorari2010,PannocchiaECC2015,PannocchiaIFAC2015}.
The closest-admissible fallback stage follows the reference tracking MPC idea of
Limon et al.\ \cite{Limon2008}.

\section{Context in the Repository}
Within this project, the target selector is used upstream of the Lyapunov MPC
tracking optimization. Its role is to convert a requested output setpoint into a
consistent equilibrium triplet
\begin{align}
    (\xs,\us,\ds),
\end{align}
where $\ds$ is taken directly from the current disturbance estimate and is not
re-optimized. The output of the target selector is then used to define the
target-centered tracking problem and terminal ingredients in the Lyapunov MPC
layer.

The code assumes the controller works in scaled deviation coordinates. This is
important: the nominal input $u_{\mathrm{nom}}$ and the steady-state target
variables are all interpreted in the same coordinate system as the MPC model.
In the current codebase, this should normally mean that
\begin{align}
    u_{\mathrm{nom}} \approx 0
\end{align}
unless a different target-centering convention has been explicitly chosen.

\section{Augmented Offset-Free Model}
The starting point is the standard disturbance-augmented linear model used in
offset-free MPC \cite{PannocchiaRawlings2003,PannocchiaBemporad2007}:
\begin{align}
    \xaug_{k+1} &= A_{\mathrm{aug}} \xaug_k + B_{\mathrm{aug}} u_k, \\
    y_k &= C_{\mathrm{aug}} \xaug_k,
\end{align}
with augmented state
\begin{align}
    \xaug = \begin{bmatrix} x \\ d \end{bmatrix},
\end{align}
where $x \in \R^{n_x}$ is the physical state and $d \in \R^{n_d}$ is the
disturbance state. In the current implementation, the augmented state ordering
is assumed to be
\begin{align}
    \xaug = \begin{bmatrix} x \\ d \end{bmatrix},
    \qquad n_d = n_y,
\end{align}
that is, the disturbance dimension is assumed equal to the output dimension.

The matrices are partitioned as
\begin{align}
    A_{\mathrm{aug}} &=
    \begin{bmatrix}
        A & B_d \\
        0 & I
    \end{bmatrix}, \\
    B_{\mathrm{aug}} &=
    \begin{bmatrix}
        B \\
        0
    \end{bmatrix}, \\
    C_{\mathrm{aug}} &=
    \begin{bmatrix}
        C & C_d
    \end{bmatrix}.
\end{align}
Equivalently, the physical model is
\begin{align}
    x_{k+1} &= A x_k + B u_k + B_d d_k, \\
    d_{k+1} &= d_k, \\
    y_k &= C x_k + C_d d_k.
\end{align}

Given the current augmented observer estimate
\begin{align}
    \xhataug = \begin{bmatrix} \hat x \\ \hat d \end{bmatrix},
\end{align}
the target selector freezes the disturbance state at
\begin{align}
    \ds = \dhat.
\end{align}
This is the standard offset-free target-calculation logic: estimate the
disturbance with the augmented observer, freeze the disturbance estimate, and
compute the equilibrium state and input consistent with that estimate
\cite{PannocchiaBemporad2007,MaederMorari2010,PannocchiaECC2015}.

\section{Steady-State Target Equations}
At steady state, the physical state and input must satisfy
\begin{align}
    (I-A)\xs - B\us - B_d \dhat = 0.
    \label{eq:ss_dyn}
\end{align}
The corresponding steady-state output is
\begin{align}
    \yexpr = C \xs + C_d \dhat.
    \label{eq:ss_output}
\end{align}

If all outputs are directly targeted, the controlled target expression is
\begin{align}
    \yctrl = \yexpr.
\end{align}
The implementation also supports an optional controlled-output selection matrix
$H$, in which case
\begin{align}
    \yctrl = H \yexpr,
\end{align}
so that the target vector $\ysp$ lives in $\R^{n_c}$ rather than $\R^{n_y}$.

\section{Constraints and Tightening}
The selector supports input bounds
\begin{align}
    u_{\min} \le \us \le u_{\max},
\end{align}
and optional tightening:
\begin{align}
    u_{\ell} &= u_{\min} + u_{\mathrm{tight}}, \\
    u_h &= u_{\max} - u_{\mathrm{tight}}.
\end{align}
These tightened input bounds are the actual bounds enforced in both target
selection stages.

Optional output bounds may also be provided:
\begin{align}
    y_{\min} \le \yexpr \le y_{\max}.
\end{align}
When output tightening is enabled, the tightened bounds become
\begin{align}
    y_{\ell} &= y_{\min} + y_{\mathrm{tight}}, \\
    y_h &= y_{\max} - y_{\mathrm{tight}}.
\end{align}

The implementation can handle the output bounds in two ways:
\begin{enumerate}[label=\alph*)]
    \item as hard inequalities, or
    \item as soft inequalities with nonnegative slack variables.
\end{enumerate}
The key design rule is that only the output inequalities may be softened. The
exact target equality in the first stage is never softened.

\section{Stage 1: Exact Offset-Free Steady-State Target}
The first stage asks whether the requested setpoint is itself an admissible
equilibrium for the disturbance-frozen model. The optimization problem is
\begin{align}
    \min_{\xs,\us} \quad
    & (\us-u_{\mathrm{nom}})^\top R_u (\us-u_{\mathrm{nom}})
    + \xs^\top Q_x \xs
    \label{eq:stage1_cost}
\end{align}
subject to
\begin{align}
    (I-A)\xs - B\us - B_d \dhat &= 0, \label{eq:stage1_dyn} \\
    \yctrl &= \ysp, \label{eq:stage1_eq} \\
    u_{\ell} \le \us &\le u_h. \label{eq:stage1_input_bounds}
\end{align}
If hard output bounds are enabled, they are also enforced:
\begin{align}
    y_{\ell} \le \yexpr \le y_h.
\end{align}

If soft output bounds are enabled, the code instead introduces
$s_{\ell}, s_h \in \R^{n_y}_{\ge 0}$ and uses
\begin{align}
    \yexpr + s_{\ell} &\ge y_{\ell}, \\
    \yexpr - s_h &\le y_h,
\end{align}
while augmenting the cost by
\begin{align}
    s_{\ell}^\top W_{\ell} s_{\ell}
    + s_h^\top W_h s_h.
\end{align}

The matrix $Q_x$ is not meant to bias the target significantly. It acts as a
small regularizer to make the state solution numerically well posed when
multiple equilibria correspond to similar inputs. The dominant Stage 1 design
principle is exact target satisfaction, not compromise between target error and
control effort.

\section{Stage 2: Closest-Admissible Fallback Target}
If Stage 1 is infeasible or numerically unacceptable, the implementation falls
back to a closest-feasible equilibrium problem. This is the correct place to
penalize target mismatch, because at this point exact tracking has already been
ruled out. The fallback problem is
\begin{align}
    \min_{\xs,\us} \quad
    & (\yctrl-\ysp)^\top T_y (\yctrl-\ysp)
    + (\us-u_{\mathrm{nom}})^\top R_u (\us-u_{\mathrm{nom}})
    + \xs^\top Q_x \xs.
    \label{eq:stage2_cost}
\end{align}
If previous target information is provided, the code adds smoothing terms
\begin{align}
    &(\xs-\xs^{\mathrm{prev}})^\top Q_{dx} (\xs-\xs^{\mathrm{prev}}) \notag\\
    &\qquad + (\us-\us^{\mathrm{prev}})^\top R_{du} (\us-\us^{\mathrm{prev}}).
    \label{eq:stage2_smoothing}
\end{align}
These regularizers are only used in the fallback stage. They are intentionally
excluded from Stage 1 so that exact target satisfaction is never weakened by
target-motion smoothing.

The fallback constraints are
\begin{align}
    (I-A)\xs - B\us - B_d \dhat &= 0, \\
    u_{\ell} \le \us &\le u_h,
\end{align}
plus the same hard or soft output-bound structure, if output bounds are present.
There is no hard equality $\yctrl = \ysp$ in Stage 2.

This formulation is aligned with the closest-admissible steady-state idea used
in constrained tracking MPC \cite{Limon2008}: if the requested target cannot be
exactly achieved, the controller should move toward the nearest admissible
steady state rather than forcing an infeasible reference.

\section{Numerical Acceptance Logic}
The code does not accept a target solely because a CVXPY solver returns
\texttt{optimal} or \texttt{optimal\_inaccurate}. Instead, it validates the
solution using explicit residual checks.

\subsection{Dynamic Residual}
For either stage, the dynamic residual is
\begin{align}
    r_{\mathrm{dyn}} = (I-A)\xs - B\us - B_d \dhat.
\end{align}
The implementation computes
\begin{align}
    \norminf{r_{\mathrm{dyn}}}.
\end{align}

\subsection{Target Equality Residual}
For Stage 1, the exact target residual is
\begin{align}
    r_{\mathrm{eq}} = \yctrl - \ysp,
\end{align}
with infinity norm
\begin{align}
    \norminf{r_{\mathrm{eq}}}.
\end{align}

\subsection{Tolerance Policy}
The implementation uses the following acceptance thresholds:
\begin{align}
    \norminf{r_{\mathrm{dyn}}} &\le 10^{-6}
    && \text{if the status is \texttt{optimal}}, \\
    \norminf{r_{\mathrm{dyn}}} &\le 10^{-5}
    && \text{if the status is \texttt{optimal\_inaccurate}}.
\end{align}
For Stage 1 the same threshold policy is applied to
\begin{align}
    \norminf{r_{\mathrm{eq}}}.
\end{align}
If hard output bounds are active, the implementation also checks the numerical
bound violation magnitude and rejects the candidate if the hard constraints are
violated beyond tolerance.

This acceptance discipline is important in practice, particularly when fallback
solves or soft-bound formulations are passed to first-order or inexact solvers.
The target selector must return a true equilibrium, not just a nominally solved
optimization variable tuple.

\section{Returned Values and Debug Information}
On success, the function returns
\begin{align}
    (\xs,\us,\ds),
    \qquad \ds = \dhat.
\end{align}
It also constructs the augmented target
\begin{align}
    x_s^{\mathrm{aug}} =
    \begin{bmatrix}
        \xs \\
        \dhat
    \end{bmatrix},
\end{align}
which is useful for the Lyapunov tracking formulation.

The debug dictionary is intentionally richer than the basic return tuple. It
includes, among other fields:
\begin{itemize}[leftmargin=2em]
    \item \texttt{solve\_stage}, indicating whether the accepted solution came
    from the exact or fallback stage,
    \item \texttt{target\_error}, \texttt{target\_error\_inf}, and
    \texttt{target\_error\_norm},
    \item \texttt{target\_eq\_residual\_inf},
    \item \texttt{dyn\_residual\_inf},
    \item \texttt{bound\_violation\_inf},
    \item \texttt{x\_s\_aug},
    \item soft-bound slack values when applicable,
    \item stage-specific status and reject-reason information.
\end{itemize}

For backward compatibility with existing diagnostics, the implementation also
defines
\begin{align}
    \texttt{target\_slack} &:= \texttt{target\_error}, \\
    \texttt{target\_slack\_inf} &:= \texttt{target\_error\_inf}.
\end{align}
These names are retained even though the new implementation no longer uses a
single explicit output-tracking slack variable in the main exact-target solve.

\section{Why the Two-Stage Structure Matters}
The main conceptual reason for the redesign is that the target selector should
not blend exact offset-free equilibrium calculation and infeasibility handling
into one weighted compromise problem. Those are two distinct tasks:
\begin{enumerate}[label=\arabic*.]
    \item If the requested setpoint is feasible for the disturbance-frozen model,
    the target selector should return the exact equilibrium corresponding to that
    setpoint.
    \item Only if the setpoint is not feasible should the target selector solve
    a nearest-admissible target problem.
\end{enumerate}

This distinction is especially important in the present repository because the
Lyapunov tracking MPC uses the selected steady-state target as the center for
tracking errors, terminal costs, and terminal-set reasoning. If the target
selector prematurely softens the target even when exact tracking is feasible,
then the Lyapunov layer is centered around an artificially biased equilibrium.

\section{Mapping Between Code Arguments and Mathematical Objects}
\begin{center}
\begin{tabular}{>{\ttfamily}p{0.31\textwidth} p{0.6\textwidth}}
\toprule
Python argument or field & Mathematical role \\
\midrule
A\_aug, B\_aug, C\_aug & Augmented offset-free prediction model \\
xhat\_aug & Current augmented estimate $[\hat x;\hat d]$ \\
y\_sp & Requested steady-state target $\ysp$ \\
u\_min, u\_max & Untightened input limits \\
u\_tight, y\_tight & Tightening vectors for target admissibility \\
u\_nom & Input-centering reference $u_{\mathrm{nom}}$ \\
Ty\_diag & Diagonal of the fallback target-mismatch weight $T_y$ \\
Ru\_diag & Diagonal of the input-centering weight $R_u$ \\
Qx\_diag & Diagonal of the state regularization matrix $Q_x$ \\
Qdx\_diag, Rdu\_diag & Previous-target regularization weights $Q_{dx}, R_{du}$ \\
y\_min, y\_max & Optional output constraints \\
Wy\_low\_diag, Wy\_high\_diag & Soft output-bound penalties $W_{\ell}, W_h$ \\
H & Optional output-selection matrix for controlled outputs \\
solve\_stage & Indicator of whether Stage 1 or Stage 2 produced the accepted target \\
\bottomrule
\end{tabular}
\end{center}

\section{Recommended Method Description for a Thesis or Paper}
A concise method description consistent with the implemented selector is:

\begin{quote}
The steady-state target selector is based on a standard offset-free MPC
formulation using a disturbance-augmented linear prediction model and
observer-based disturbance estimation. At each control step, the current
disturbance estimate is frozen and an exact steady-state target problem is
solved for the equilibrium state and input consistent with the requested
setpoint. If the requested setpoint is infeasible under the disturbance-frozen
model and the imposed constraints, the selector computes the closest admissible
steady-state target instead. This yields a target calculation that is exact when
feasible and only softened when infeasibility makes exact tracking impossible.
\end{quote}

\section{Compilation Notes}
This report is written as a self-contained \LaTeX{} file with an internal
bibliography, so it can be compiled with plain \texttt{pdflatex} without
\texttt{bibtex} or \texttt{biber}. On a machine with a \LaTeX{} installation,
run:
\begin{verbatim}
cd report
pdflatex -interaction=nonstopmode target_selector_report.tex
pdflatex -interaction=nonstopmode target_selector_report.tex
\end{verbatim}
The second pass resolves cross-references and citation numbering cleanly.

\begin{thebibliography}{99}

\bibitem{PannocchiaRawlings2003}
G.~Pannocchia and J.~B.~Rawlings,
``Disturbance models for offset-free model-predictive control,''
\emph{AIChE Journal}, 2003.

\bibitem{PannocchiaBemporad2007}
G.~Pannocchia and A.~Bemporad,
``Combined design of disturbance model and observer for offset-free model predictive control,''
\emph{IEEE Transactions on Automatic Control}, 2007.

\bibitem{MaederMorari2010}
U.~Maeder and M.~Morari,
``Offset-free reference tracking with model predictive control,''
\emph{Automatica}, 2010.

\bibitem{PannocchiaECC2015}
G.~Pannocchia,
``Offset-free tracking MPC: A tutorial review and comparison of different formulations,''
in \emph{Proceedings of the European Control Conference}, 2015.

\bibitem{PannocchiaIFAC2015}
G.~Pannocchia, M.~Gabiccini, and A.~Artoni,
``Offset-free MPC explained: novelties, subtleties, and applications,''
\emph{IFAC-PapersOnLine}, 2015.

\bibitem{Limon2008}
D.~Limon, I.~Alvarado, T.~Alamo, and E.~F.~Camacho,
``MPC for tracking piecewise constant references for constrained linear systems,''
\emph{Automatica}, 2008.

\end{thebibliography}

\end{document}
